\section{序論と主結果}

半直交部分圏の変形が、もとの周囲の圏の変形から常に誘導されるとは限らない。
本稿ではこの問題を、滑らかな重み付き射影スタック
\[
 X=\mathbb P(1,1,1,3)
\]
の六項例外列と、その四項部分列について具体的に解く。
係数体は全篇を通して $k=\mathbb C$ とする。
重み付き多項式環と二つの有限代数を
\[
 R=k[x_1,x_2,x_3,w],\quad \deg x_i=1,\quad\deg w=3,
 \qquad \Lambda=R_{(0,1,2,3,4,5)},\quad B=R_{(0,1,3,4)}
\]
で定める。次数切り出しの対象 $i<j$ 間の射空間は対応する次数差の $R$ の成分であり、
合成は多項式積である。六項の番号を保つ冪等元
$e=e_0+e_1+e_3+e_4$ に対して $B=e\Lambda e$ となる。
$B$ が本稿でいうVeronese--Kronecker代数である。

Auroux--Katzarkov--Orlovの定理\cite[Theorem 2.12]{AKO}により、
$(\mathcal O_X,\ldots,\mathcal O_X(5))$ はfull strong exceptional collectionであり、
\[
 \Perf(\Lambda)\simeq\Perf(X),\qquad
 \Perf(B)\hookrightarrow\Perf(\Lambda)
\]
は許容部分圏の埋め込みである。
本稿の補完問題は、特殊ファイバー上のこの埋め込みを固定し、
与えられた $B$ の平坦変形を $\Lambda$ の変形のcornerとして実現できるか、という問題である。
頂点冪等元と特殊ファイバーの同定を固定する。
別の例外列や、この特殊ファイバーの埋め込みと結び付かない別の埋め込みについては主張しない。

\subsection{接空間の一致と三次コホモロジーの差}

制限写像を $\rho^i:HH^i(\Lambda)\to HH^i(B)$ と書く。
本稿では次の値を得る。
\[
\begin{array}{c|rrr}
 &HH^1&HH^2&HH^3\\\hline
 \Lambda&19&53&41\\
 B&19&53&39\\
 \operatorname{rank}\rho&19&53&1
\end{array}
\]
特に $\rho^1,\rho^2$ は同型である。
しかし $\ker\rho^3$ は四十次元であり、周囲の三次情報のほとんどが四項側で失われる。
非自明慣性成分を除いた部分には自然な同定
\[
 HH^3(X)_{\rm untw}=H^0(X,\mathcal O_X(6))
 =\operatorname{Sym}^6\langle x_1,x_2,x_3\rangle
 \oplus w\operatorname{Sym}^3\langle x_1,x_2,x_3\rangle\oplus kw^2
\]
があり、この部分の制限の核は最初の二項、すなわち三十八次元である。
これはstacky点 $P=(0:0:0:1)$ で消える反標準切断の空間でもある。

この三十八次元部分は、$HH^2(X)$ の二つの元のSchouten括弧が張る空間に正確に一致する。
したがって、周囲の一次変形に現れる全ての二次障害は四項側で消える。
三十八次元の二次障害の線形包と、四十次元の全制限核は区別する必要がある。

Kuznetsovのnormal Hochschild cohomologyは、半直交分解から生じる制限と変形空間の比較を
一般的に説明する\cite[\S4.2]{Kuznetsov}。
本稿は、二次コホモロジーが一致しても三次で単射性が失われる場合の、具体的な補完軌跡を与える。

\subsection{全次数の変形空間}

$V=H^0(X,\wedge^2T_X)$ と置くと $\dim V=53$ である。
そのJacobi式 $[\pi,\pi]/2=0$ の独立な係数は三十八個の同次二次式
$q_1,\ldots,q_{38}$ となる。
本稿の中心的な結論は、適切な形式座標とcompatibleなminiversal族を選ぶと、corner制限が
\begin{equation}\label{eq:intro-hulls-v31}
 \operatorname{Spf}\frac{k[[u_1,\ldots,u_{53}]]}{(q_1,\ldots,q_{38})}
 \ \hookrightarrow\ \operatorname{Spf}k[[u_1,\ldots,u_{53}]]
\end{equation}
と表示されることである。
左が六項側、右が四項側のminiversal基底である。
これは非標準的なminiversal基底の表示であり、自己同型を持つ変形関手のfine表現ではない。
また三十八は最小二次生成元の数であって、余次元でも正則列の長さでもない。

証明では、重み零多重ベクトルに発散零代表を選び、アフィン空間のKontsevich量子化
\cite{Kontsevich}を全Jacobi軌跡上で同時に適用する。
重み付き発散分解はTang--Wang--Zhang\cite{TWW}、通常の射影空間における対応は
Pym\cite[\S2.2]{Pym}の研究に位置付けられる。
これらの分解や量子化定理自体を新しいものとは扱わない。
ここで必要なのは、それを六項有限代数へ切り出し、二次障害の完全な計算と
Nakayamaの補題を組み合わせて、全次数のminiversal基底を認定することである。

式\eqref{eq:intro-hulls-v31}の重要な帰結として、任意のPoisson一次方向には
補完可能な全形式VK変形が少なくとも一つ存在する。
一方、指定された全形式VK変形の補完可能性は、その一次方向だけでは決まらない。

\subsection{節点による正確な補完条件}

本稿で構成する二変数の階数五十六の平坦族 $B_{s,t}$ については、任意の局所Artin基底 $T$ で
\[
 B_{s,t}\otimes T\text{ が標準六項へ補完可能}
 \quad\Longleftrightarrow\quad st=0\text{ in }T
\]
が成立する。必要性は、特定のOre表示を仮定せず、六項有限代数の全ての変形に対して示す。
証明には $\rho^1$ の全射性と $\rho^2$ の単射性から得られるcorner同型の持上げを使う。

特に $r\ge2$ に対する形式弧 $B_{\hbar,\hbar^r}$ は、非零Poisson一次方向を持つが、
$\hbar^{r+1}$ を法としては補完でき、$\hbar^{r+2}$ を法として初めて補完できなくなる。
従って、全ての形式弧に一様に適用できる固定有限次数のjet判定は存在しない。
この事実は、式 $st=0$ のような有限本の形式方程式による記述と両立する。


\subsection{補完空間の幾何との接続}
Poissonスキーム自体の成分と局所環は、独立した別稿\cite{VKPoisson32}で扱う。
そこでは八つの相異なる成分を認定し、そのうち一つの十九次元成分の一般点に
長さ六の非被約構造を見いだす。本稿のhull同定により、この幾何は六項補完空間にも移る。
これにより、四項側の滑らかさに対して、六項側が非純次元かつ非被約であることが分かる。
非被約構造を捨てたPoisson構造の集合だけでは、Artin基底上の補完問題を記述できない。
本稿では変形理論と補完条件の証明を完結させ、成分ごとの詳細な幾何は上記別稿へ分ける。

\subsection{規約と計算の役割}

多重ベクトルのHKR写像は $p!$ で割った反対称化を用いる。
従って変形乗法 $m+\varepsilon\mu_1$ の半古典括弧は
$\mu_1(f,g)-\mu_1(g,f)$ である。
計算付録の評価行列は整数を保つため階乗を省くが、その階数と像は変わらない。
具体的な二変数族では、この規約により半古典括弧に係数二が現れる。

有限計算は全て有理数の厳密演算で再現できる。
それは接空間、制限写像、独立な障害式を検査するものであり、
全次数の量子化、Artin基底上の必要十分性、形式基底の同型は本文の構造的証明による。
